% =====================================================================
% World Model Sidewalk: barebones paper skeleton
% Format: IEEE conference (ICRA/IROS/RA-L style). Compile with:
%   latexmk -pdf main.tex     (or pdflatex -> bibtex -> pdflatex x2)
% Status: DRAFT. Sections marked \todo{} are placeholders.
% =====================================================================
\documentclass[conference]{IEEEtran}

\usepackage[utf8]{inputenc}
\usepackage[T1]{fontenc}
\usepackage{amsmath,amssymb}
\usepackage{graphicx}
\usepackage{booktabs}
\usepackage{xcolor}
\usepackage{cite}
\usepackage[hidelinks]{hyperref}

\newcommand{\todo}[1]{\textcolor{red}{[TODO: #1]}}
\newcommand{\vect}[1]{\mathbf{#1}}

\begin{document}

\title{Planning in Latent Space for Legged Navigation:\\
Action-Conditioned JEPA World Models for Quadruped Local Planning\\
on Sidewalk and Construction-Like Terrain}

\author{\IEEEauthorblockN{Elyas Yassin}
\IEEEauthorblockA{University of Colorado Denver\\
\todo{email / co-authors / advisor}}}

\maketitle

% ---------------------------------------------------------------------
\begin{abstract}
Legged robots deployed on sidewalks and active work sites must handle
terrain whose effect on motion is hard to model by hand: curbs, slopes,
low-friction patches, loose debris, and layouts that change faster than
maps are updated. Hand-written kinematic models assume commanded and
realized motion agree, which fails exactly where these robots are most
useful. We study whether an action-conditioned latent world model,
built on frozen V-JEPA~2 visual features fused with proprioception, can
serve as the dynamics model for a short-horizon local planner on a
Unitree Go2 quadruped, while a conventional A* planner handles
route-level navigation and a pretrained locomotion policy executes
velocity commands. The model predicts future latent and physical states
without reconstructing images, so candidate command sequences can be
scored at planning rates. We evaluate two hypotheses in simulation:
(H1) visual plus proprioceptive features predict navigation outcomes on
variable terrain better than proprioception alone, and (H2) learned
model predictive control improves safety and success relative to
kinematic and reactive local planners. Because the local planner
optimizes an explicit cost at run time, we also test (H3) whether
spoken instructions, grounded in a scene graph and translated into
structured goals, keep-out regions, and speed limits, can re-task the
robot mid-mission without retraining and without weakening safety
constraints. Evaluation uses a terrain
difficulty ladder that trains on sidewalk scenes and tests on held-out
construction-like terrain. Preliminary results on flat ground show that
a learned state-only model predicts motion more accurately than command
kinematics, yet yields no navigation advantage, because failures are
dominated by locomotion-policy deadband and planner cost tuning. These
findings motivate the terrain ladder as a test of where, if anywhere,
learned latent dynamics begin to pay off.
\todo{update with H1/H2 results once Phase E runs}
\end{abstract}

\begin{IEEEkeywords}
world models, legged robots, model predictive control, visual navigation,
self-supervised learning, traversability
\end{IEEEkeywords}

% ---------------------------------------------------------------------
\section{Introduction}
\todo{Motivation: sidewalk delivery/inspection and construction-site
documentation; why hand-modeled dynamics break on unstructured terrain.}

\todo{Gap: action-conditioned latent world models
\cite{assran2025vjepa2,zhou2024dinowm,zholus2026robojepa} plan well on
manipulators and mazes but have not been used as the test-time dynamics
model for legged local navigation.}

\todo{Contributions (draft):}
\begin{enumerate}
  \item A hierarchical navigation stack that uses an action-conditioned
        JEPA-style latent world model as the dynamics model inside a
        receding-horizon planner on a quadruped, beneath an A* global
        planner and above a pretrained locomotion policy.
  \item A controlled comparison of state-only and visual-latent dynamics
        models (H1) and of learned, kinematic, and reactive local planners
        (H2) on identical seeded scenes with paired statistics.
  \item A terrain difficulty ladder, from flat ground to construction-like
        terrain with unmapped layout changes, that locates the conditions
        under which learned latent dynamics outperform hand-written models.
  \item A human-in-the-loop interface in which spoken commands, grounded
        in a scene graph, become structured constraints on the global and
        local planners, enabling mid-mission re-tasking that is always
        gated by the safety filter.
  \item A diagnosis of execution-layer confounds (locomotion-policy
        deadband, clearance cost tuning) that can mask or mimic planner
        differences in legged-navigation benchmarks.
\end{enumerate}

% ---------------------------------------------------------------------
\section{Related Work}

\subsection{Latent world models for planning}
Joint-embedding predictive architectures (JEPAs) learn by predicting
representations rather than pixels \cite{lecun2022path,assran2023ijepa}.
V-JEPA~2 pre-trains a video encoder on over one million hours of
internet video and post-trains an action-conditioned predictor
(V-JEPA~2-AC) on under 62~hours of robot video, enabling zero-shot
pick-and-place on Franka arms by planning toward image goals
\cite{assran2025vjepa2}. DINO-WM predicts future DINOv2 patch features
from offline trajectories and plans by optimizing action sequences
toward goal features in mazes and manipulation tasks
\cite{zhou2024dinowm}. PLDM compares JEPA latent-dynamics planning with
model-free offline RL on navigation tasks and finds planning more data
efficient and better at generalizing to unseen layouts
\cite{sobal2025pldm}. RoboJEPA scales a JEPA predictor to 8B parameters
across 12 embodiments and shows that imagination error follows a power
law in compute and predicts downstream planning success
\cite{zholus2026robojepa}. TD-MPC2 performs trajectory optimization in
the latent space of a decoder-free model learned jointly with a value
function \cite{hansen2024tdmpc2}, and PETS established sampling-based
MPC with learned probabilistic dynamics \cite{chua2018pets}.
\emph{Open:} these methods are demonstrated on arms, simulated mazes,
and benchmark control suites. None uses a latent world model as the
test-time dynamics model for a legged robot navigating outdoor terrain,
and none feeds proprioception (body tilt, contacts, realized velocity)
into the latent rollout.

\subsection{Generative world models for navigation}
Navigation World Models (NWM) train a 1B-parameter conditional diffusion
transformer to generate future egocentric frames and rank candidate
trajectories \cite{bar2025nwm}. DreamerV3 learns a recurrent
state-space model with reconstruction for actor-critic learning
\cite{hafner2023dreamerv3}, and DreamerNav applies Dreamer-style world
models to navigation in dynamic indoor environments
\cite{shanks2025dreamernav}. WAVE-Go couples world-model action
prediction with interruptible, risk-budgeted execution for wheel-legged
image-goal navigation \cite{li2026wavego}.
\emph{Open:} pixel generation is expensive to sample inside an MPC loop,
and these models do not represent the execution limits of the
underlying locomotion controller.

\subsection{World models and perception for legged locomotion}
DayDreamer learns Dreamer world models online on physical robots,
including a quadruped \cite{wu2022daydreamer}. WMP trains a world model
in simulation and conditions a locomotion policy on its latent state,
improving traversability over teacher-student baselines
\cite{lai2024wmp}. Egocentric-vision locomotion policies cross stairs,
curbs, and gaps \cite{agarwal2022egocentric}, and coupling vision with
proprioception improves legged navigation around obstacles
\cite{fu2022coupling}. Most closely related, Zhu et al.\ add a
JEPA-style auxiliary predictor to a recurrent RL navigation policy on a
Go2; the predictor shapes training and is discarded at inference
\cite{zhu2026predictive}.
\emph{Open:} in these works the world model shapes a policy or a gait
rather than serving as an explicit model queried at test time to score
navigation command sequences.

\subsection{Visual traversability learning}
Wild Visual Navigation learns traversability online from pre-trained
visual features with proprioceptive self-supervision
\cite{mattamala2024wvn}, and GrandTour provides large-scale multi-modal
legged-robot data across urban areas, forests, and demolished buildings
\cite{tuna2026grandtour}.
\emph{Open:} traversability estimators score instantaneous terrain
appearance for a geometric planner; they do not predict how the robot's
state evolves under a proposed command sequence.

\subsection{Navigation foundation models}
ViNT \cite{shah2023vint} and NoMaD \cite{sridhar2024nomad} learn
goal-conditioned navigation policies from hundreds of hours of
multi-robot data.
\emph{Open:} they map observations to actions without an explicit
predictive model, which makes it hard to inject new constraints at test
time or attribute failures to perception, planning, or execution.

\subsection{Scene graphs and language-conditioned planning}
3D scene graphs represent environments as layered graphs of objects,
places, and rooms \cite{armeni2019scenegraph}, and can be built in real
time from sensor data \cite{hughes2022hydra} or with open-vocabulary
labels from foundation models \cite{gu2024conceptgraphs}. SayPlan grounds
LLM task planning in a 3D scene graph and delegates paths to a classical
planner \cite{rana2023sayplan}; VLMaps fuses visual-language features
into a metric map so language goals can be localized for navigation
\cite{huang2023vlmaps}. On the control side, Code as Policies has LLMs
write robot programs over primitive skills \cite{liang2023cap}, and
Language to Rewards has LLMs specify reward parameters that a real-time
MPC optimizer turns into behavior, including on a simulated quadruped,
with interactive corrections from the user \cite{yu2023l2r}. NARRATE has
an LLM write MPC objectives and constraints as mathematical expressions
for manipulation \cite{ismail2024narrate}. Closest to our interface,
``Hey Robot!'' uses a vision-language model to generate the cost
function and reconfigure the parameters of a navigation MPC from user
instructions on a ground robot \cite{martinezbaselga2024heyrobot}. For
legged robots, Kulkarni et al.\ cage a vision-language navigation
planner with an MPPI planner under barrier-function constraints for
provable collision avoidance \cite{kulkarni2026caging}.
\emph{Open:} translating language into MPC costs is established. In
these systems the MPC relies on analytic or simulator dynamics and
geometric maps; none couples language-specified constraints with a
learned, terrain-aware dynamics model, and grounding of spatial
references in a scene graph and mid-mission re-tasking by voice are
not evaluated together. \todo{Verify whether Hey Robot supports
mid-mission re-configuration; adjust this claim.}

\subsection{Positioning}
This work sits at the intersection of the first and third threads: a
non-generative, action-conditioned latent world model used for
receding-horizon planning on a legged robot, evaluated against
kinematic and reactive planners with explicit control of execution-layer
confounds. A secondary contribution connects this planner to spoken
instructions. Language-to-MPC interfaces exist
\cite{martinezbaselga2024heyrobot,ismail2024narrate}; what we add is
coupling them with a learned terrain-aware dynamics model on a legged
robot, scene-graph grounding, and safety-gated mid-mission re-tasking,
none of which an end-to-end navigation policy supports without
retraining. \todo{Confirm with a targeted search that no prior work plans
with a JEPA latent model on a legged robot; record the search bound.}

% ---------------------------------------------------------------------
\section{Proposed Methodology}

\subsection{Problem formulation}
The robot receives a goal and a 2D occupancy map that may be outdated.
At control step $t$ (10~Hz) the local planner outputs a body-frame
velocity command $\vect{a}_t = (v_x, v_y, \omega_z)$, which a fixed
pretrained locomotion policy $\pi_{\text{loco}}$ converts to joint
targets. We do not learn gait control. The observation
$\vect{o}_t = (\vect{s}_t, I_{t-T:t})$ contains a proprioceptive state
$\vect{s}_t$ (base pose, realized velocity, IMU, and contacts where
available) and an RGB context window $I_{t-T:t}$ from a body-mounted
camera. The objective is to reach the goal without collisions or falls,
minimizing time and hazard exposure.

\subsection{System architecture}
The stack has three levels \todo{Fig.~1: architecture diagram}:
\begin{enumerate}
  \item \textbf{Global planner:} A* \cite{hart1968astar} on the occupancy
        grid produces waypoints and triggers a reroute when the local
        planner reports a blocked path.
  \item \textbf{Local planner:} a receding-horizon planner (1--3~s)
        optimizes command sequences against a dynamics model and
        executes the first command.
  \item \textbf{Execution:} $\pi_{\text{loco}}$ tracks the command in
        Isaac Lab \cite{nvidia2025isaaclab}, using a policy trained with
        massively parallel RL \cite{rudin2021walk}.
\end{enumerate}

\subsection{Dynamics models}
\paragraph{State-only baseline}
An MLP or GRU predicts the change in local-frame state,
$\Delta\hat{\vect{s}}_{t+1} = f_\theta(\vect{s}_{t-h:t}, \vect{a}_{t-h:t})$,
from a short history ($h{=}4$, 0.4~s), trained with a multi-step rollout
loss. Heading is represented by $(\sin\psi, \cos\psi)$ to avoid
wrap-around.

\paragraph{Visual latent model}
A frozen V-JEPA~2 encoder $E_v$ \cite{assran2025vjepa2} embeds the RGB
window, and a small encoder $E_p$ embeds proprioception:
\begin{equation}
  \vect{z}_t = g_\phi\big(E_v(I_{t-T:t}),\, E_p(\vect{s}_t)\big).
\end{equation}
An action-conditioned predictor $P_\psi$ rolls the latent forward given
the command sequence and its durations:
\begin{equation}
  \hat{\vect{z}}_{t+k} = P_\psi\big(\vect{z}_t, \vect{a}_{t:t+k-1}, \Delta t_{t:t+k-1}\big).
\end{equation}
Training targets come from the frozen encoder with a stop-gradient,
which avoids representation collapse without extra regularization:
\begin{multline}
  \mathcal{L} = \lambda_z\, d\big(\hat{\vect{z}}^{v}_{t+k},\, \mathrm{sg}[E_v(I_{t+k-T:t+k})]\big) \\
  + \lambda_s\, \lVert h_s(\hat{\vect{z}}_{t+k}) - \vect{s}_{t+k} \rVert^2
  + \lambda_r\, \mathcal{L}_{\text{risk}},
\end{multline}
where $h_s$ decodes physical state and $\mathcal{L}_{\text{risk}}$ is
included only where simulator-derived labels (collision, fall, slip)
are valid. Linear probes test whether task variables (distance to
obstacle, realized speed, traversability) are decodable from
$\vect{z}_t$; low latent loss alone is not taken as evidence of
planning utility.

\subsection{Local planner}
We use the cross-entropy method (CEM) \cite{rubinstein1999cem} over
bounded command sequences of horizon $H$. Each candidate is rolled out
through the dynamics model and scored by
\begin{equation}
  J = w_g\, d_{\text{wp}} + w_c\, c_{\text{clear}} + w_r\, \hat{r}
    + w_u\, \lVert \Delta \vect{a} \rVert^2 + w_\sigma\, \hat{\sigma},
\end{equation}
covering waypoint distance, a clearance penalty, predicted risk (only
when validated), command smoothness, and calibrated model uncertainty.
Two execution-aware components follow from preliminary experiments:
\begin{itemize}
  \item \textbf{Deadband projection:} sampled planar commands in the
        policy's measured deadband are mapped to a stop or to the
        smallest command the policy reliably realizes.
  \item \textbf{Safety filter:} a geometric hard constraint on obstacle
        proximity overrides any learned cost; a control barrier function
        formulation \cite{ames2017cbf} is a candidate extension.
\end{itemize}

\subsection{Human-in-the-loop: language-specified mission constraints}
An existing voice-to-voice retrieval-augmented generation (RAG) pipeline
transcribes operator speech, interprets it, and replies by voice.
\todo{Describe the current pipeline: ASR, LLM, retrieval corpus, TTS,
latency.} It currently issues raw velocity commands (e.g., ``forward
2~m''). We restructure its output into a typed \emph{mission
specification} that the planners consume, so that language never
bypasses the local planner or safety filter:
\begin{center}
\begin{tabular}{@{}ll@{}}
\toprule
Mission-spec field & Consumed by \\
\midrule
\texttt{goal}: scene-graph node or pose & global planner \\
\texttt{via}: ordered waypoints / nodes & global planner \\
\texttt{keep\_out}: regions around nodes & global + local cost \\
\texttt{speed\_limit}: m/s, optionally per region & local bounds \\
\texttt{motion}: relative move (``forward 2~m'') & local planner goal \\
\texttt{control}: pause / resume / abort & executive \\
\bottomrule
\end{tabular}
\end{center}
\paragraph{Scene-graph grounding}
A scene graph with object, place, and region layers
\cite{armeni2019scenegraph,hughes2022hydra} resolves spatial references
(``the dumpster by the north gate'') to nodes and poses. In simulation
the graph is generated from scene metadata; building it from perception
\cite{gu2024conceptgraphs} is out of scope.
\paragraph{Compilation into the planners}
\texttt{goal}/\texttt{via} replace the A* target; \texttt{keep\_out}
adds cells to the occupancy grid and a soft penalty $w_k\,c_{\text{keep}}$
to $J$; \texttt{speed\_limit} tightens the CEM sampling bounds;
\texttt{motion} becomes a relative local goal planned by the MPC rather
than an open-loop velocity. Language-derived terms can only add costs
or tighten bounds: they cannot relax the safety filter. Commands that
are infeasible or conflict with safety (e.g., ``go straight'' into a
blocked region) are refused with a spoken explanation.
\paragraph{Mid-mission re-tasking}
A new spec preempts the current one at the next planning step; the
local planner replans from the current state without retraining.

\subsection{Terrain difficulty ladder}
Training data come from sidewalk-like scenes. Evaluation climbs a ladder
of increasingly hard-to-model terrain, with the upper rungs held out
entirely:
\begin{center}
\begin{tabular}{@{}cll@{}}
\toprule
Level & Conditions & Split \\
\midrule
L0 & Flat ground, cylinder obstacles & train / test \\
L1 & Curbs, ramps, slopes & train / test \\
L2 & Friction patches (wet, icy, gravel proxies) & train / test \\
L3 & Debris and rubble height-fields & held out \\
L4 & L3 plus unmapped layout changes & held out \\
\bottomrule
\end{tabular}
\end{center}
L3 and L4 approximate construction-site conditions within what the
physics engine can simulate faithfully (rigid height-fields, friction).
Deformable ground such as mud or sand is out of scope. \todo{L3--L4 may
require a rough-terrain locomotion policy; record which policy is used
per level.}

\subsection{Data collection}
Episodes are collected in Isaac Lab with a body-mounted RGB camera,
randomized exploration commands, and seeded scene generation per
family. Each transition stores aligned $(\vect{o}_t, \vect{a}_t,
\vect{o}_{t+1})$ at 10~Hz with timestamps and availability masks.
Splits are scene-disjoint, and test scenes are never used for tuning.

\subsection{Evaluation protocol}
\paragraph{H1 (prediction)}
Open-loop position and heading error at 0.5/1/2/3~s, and probe accuracy
for task variables, for state-only versus visual-latent models, per
ladder level. Comparisons are paired per window with bootstrap
confidence intervals over episodes.

\paragraph{H2 (navigation)}
Four local planners share the same global planner, safety filter, and
seeded scenes: (A) reactive pure pursuit with potential fields,
(B) kinematic CEM-MPC, (C) state-only learned CEM-MPC, and
(D) visual-latent CEM-MPC. Metrics are success rate, collisions, falls,
time to goal, minimum clearance, safety-filter interventions, and planner
latency (p50/p95). Paired success is compared with exact McNemar tests.
All baselines receive equivalent tuning on validation scenes.

\paragraph{H3 (language re-tasking)}
Scripted operator utterances are injected at random mission times
(re-goal, add keep-out, speed limit, relative move, pause, and
deliberately unsafe commands). Metrics: parse and grounding accuracy,
constraint satisfaction (keep-out violations, speed-limit exceedance),
success after re-tasking, latency from end of utterance to behavior
change, and correct refusal of unsafe commands. The comparison point is
the current raw-velocity interface executed open loop. Scripted
utterances avoid a human-subjects study; a user study is future work.

\subsection{Preliminary findings}
In simulation on flat ground (L0), a state-only MLP predicts Go2 motion
more accurately than command kinematics at every horizon (paired, all
test episodes), but tripling the data gives no reliable gain at 1--3~s.
Across 51 shared navigation scenes, reactive, kinematic-MPC, and
learned-MPC planners were statistically indistinguishable; every failure
was a freeze in front of a narrow gap, caused by the locomotion policy
ignoring small planar commands and by a conservative clearance margin.
\todo{Replace with confirmatory run on fresh scenes and tuned reactive
baseline; add table.}

% ---------------------------------------------------------------------
\section{Experiments}
\todo{Results per ladder level for H1 and H2.}

\section{Discussion and Limitations}
\todo{Sim-only evaluation; pretrained locomotion policy as a fixed
confound; rigid-terrain physics; no real-robot actuation in this work.}

\section{Conclusion}
\todo{}

% ---------------------------------------------------------------------
\bibliographystyle{IEEEtran}
\bibliography{references}

\end{document}
